\documentclass[11pt]{article}

\usepackage[a4paper,margin=1in]{geometry}
\usepackage{amsmath}
\usepackage{booktabs}
\usepackage{longtable}
\usepackage{array}
\usepackage{enumitem}
\usepackage{xcolor}
\usepackage{hyperref}
\usepackage{listings}

\hypersetup{
  colorlinks=true,
  linkcolor=blue,
  urlcolor=blue
}

\lstset{
  basicstyle=\ttfamily\small,
  frame=single,
  breaklines=true,
  columns=fullflexible,
  keepspaces=true,
  showstringspaces=false
}

\newcommand{\ms}{\,\mathrm{ms}}
\newcommand{\gb}{\,\mathrm{GB}}
\newcommand{\mb}{\,\mathrm{MB}}
\newcommand{\code}[1]{\texttt{#1}}

\title{Uniform Backward Pass Optimization Summary}
\author{gpurec profiling notes}
\date{May 2026}

\begin{document}
\maketitle

\section{Scope and headline results}

This note summarizes the optimization work documented in
\code{docs/uniform-backward-fp32-fused-profile.md} and
\code{docs/uniform-backward-50tree-wave2-profile.md}.  The focus is the CUDA
uniform backward path.  The tables use the benchmark timings recorded in those
documents; when a gain is called ``incremental'', it is the adjacent comparison
from the same benchmark run, not a mathematically additive global attribution.

The high-level result is:

\begin{center}
\begin{tabular}{lrrr}
\toprule
Benchmark & Early baseline & Final documented default & Gain \\
\midrule
10 families & $73.252\ms$ & $37.204\ms$ & $36.048\ms$ faster, $49.2\%$ \\
50 families & $230.795\ms$ & $121.016\ms$ & $109.779\ms$ faster, $47.6\%$ \\
50-family peak memory & $16.806\gb$ & about $10.308\gb$ & $6.498\gb$ lower, $38.7\%$ \\
\bottomrule
\end{tabular}
\end{center}

Most wins came from reducing memory traffic, removing scratch tensors, and
removing host synchronization points.  The hot kernels were not limited by
matrix multiply throughput.  Nsight Compute repeatedly showed memory bandwidth,
long-scoreboard stalls, barriers, and register pressure as the relevant
constraints.

\section{Uniform backward dataflow}

Each backward wave has three important families of work:

\begin{enumerate}[leftmargin=*]
\item The self-loop or wave kernel, \code{\_wave\_backward\_uniform\_kernel}.
      It applies the fixed-point Neumann VJP for the wave and accumulates
      parameter contributions.
\item The DTS split work.  The backward pass recomputes DTS forward quantities
      for split waves and then runs \code{\_dts\_cross\_backward\_accum\_kernel}
      to add direct \code{Pi} adjoints and produce the adjoints needed by the
      uniform \code{Pibar} VJP.
\item The uniform \code{Pibar} VJP tree kernel.  It propagates split-side
      \code{Pibar} adjoints through the species-tree correction.
\end{enumerate}

Before the later optimizations, a 50-family Nsight Systems capture broke the
kernel work down as follows:

\begin{center}
\begin{tabular}{lrr}
\toprule
Component & Time & Share of kernel time \\
\midrule
\code{\_wave\_backward\_uniform\_kernel} & $39.578\ms$ & $27.7\%$ \\
\code{\_uniform\_cross\_pibar\_vjp\_tree\_kernel} & $37.954\ms$ & $26.6\%$ \\
\code{\_dts\_cross\_backward\_accum\_kernel} & $29.678\ms$ & $20.8\%$ \\
\code{\_dts\_fused\_kernel} & $12.136\ms$ & $8.5\%$ \\
PyTorch float reductions & $5.498\ms$ & $3.9\%$ \\
PyTorch active-mask abs/max work & $7.910\ms$ & $5.5\%$ \\
\bottomrule
\end{tabular}
\end{center}

The optimization strategy was therefore to keep the fused Triton/CUDA path, but
remove avoidable dense intermediates, redundant passes, host scalar reads, and
scratch-buffer round trips.

\section{Optimization timeline}

\begin{longtable}{p{0.18\linewidth}p{0.19\linewidth}p{0.17\linewidth}p{0.17\linewidth}p{0.19\linewidth}}
\toprule
Change & Local comparison & 10-family gain & 50-family gain & Why it helped \\
\midrule
\endfirsthead
\toprule
Change & Local comparison & 10-family gain & 50-family gain & Why it helped \\
\midrule
\endhead

Leaf species index & $73.252 \to 68.373\ms$ and $230.795 \to 225.373\ms$ &
$4.879\ms$ & $5.422\ms$ & Avoided dense $[W,S]$ leaf tensors and their reads. \\

Read-only RHS, phase 1 & $68.373 \to 67.984\ms$ and $225.373 \to 218.372\ms$ &
$0.389\ms$ & $7.001\ms$ & Made the fused wave kernel treat the RHS as read-only, reducing snapshot/copy pressure. \\

Empty spec scratch & $67.984 \to 68.160\ms$ and $218.372 \to 214.755\ms$ &
neutral at 10 & $3.617\ms$ & Allocated scratch with \code{empty} instead of zeroing buffers that the kernel overwrites. \\

Direct global reductions & $68.160 \to 65.428\ms$ and $214.755 \to 205.489\ms$ &
$2.732\ms$ & $9.266\ms$ & For $G=1$, replaced indexed \code{scatter\_add\_} style reductions with direct sums/adds. \\

Fused wave parameter accumulation & $65.428 \to 60.614\ms$ and $205.489 \to 181.570\ms$ &
$4.814\ms$ & $23.919\ms$ & Removed six full $[W,S]$ contribution tensors and their PyTorch reductions. \\

Default \code{max\_wave\_size=32768} & $60.614 \to 61.095\ms$ and $181.570 \to 181.255\ms$ &
neutral & $0.315\ms$ plus $3.774\gb$ lower peak & Capped wave batches to keep large runs under memory limits with little speed cost. \\

Compact \code{Pibar} scratch & $61.095 \to 59.462\ms$ and $181.255 \to 174.174\ms$ &
$1.633\ms$ & $7.081\ms$ & Stored one compact coefficient instead of three scratch arrays. \\

RHS view, final per-wave clone removal & $48.993 \to 48.340\ms$ and $163.035 \to 159.055\ms$ &
$0.653\ms$ & $3.980\ms$ & Removed exactly one device-to-device copy per processed wave; 50-family D2D bytes fell from $2609.790\mb$ to $35.638\mb$. \\

Triton active-mask builder & $157.820 \to 151.909\ms$ at 50 families in the final default-flip check &
about $0.8\ms$ & $4.9$--$5.9\ms$ & Replaced PyTorch \code{abs().max()} plus compare kernels with one row-reduction Triton kernel. \\

Merged DTS S-term accumulation & $153.198 \to 148.625\ms$ at 50 families in the final forced-old comparison &
about $1.2\ms$ & $3.1$--$4.6\ms$ & Reused values already live in the DTS species loop instead of making a second species pass. \\

Device-resident DTS scalar parameters & about $149.6 \to 148.4\ms$ at 50 families &
about $2.5\ms$ & $1.0$--$1.5\ms$ & Removed 66 CUDA scalar-to-Python extractions in the 50-family split waves. \\

In-kernel DTS scalar reduction accumulation & $148.449 \to 147.563\ms$ mean at 50 families &
neutral & about $0.9\ms$ mean, median neutral & Removed small PyTorch scalar reductions for large split waves; kept vector \code{grad\_mt} accumulation disabled. \\

DTS-staged \code{Pibar} VJP & $44.802 \to 40.611\ms$ and $148.387 \to 136.886\ms$ &
$4.191\ms$ & $11.501\ms$ & DTS writes \code{u\_d} and row sums directly, so the tree VJP skips denominator construction and large reads. \\

Wave speciation gather, int32 topology, 8-warps, row-max reuse & $40.050 \to 37.204\ms$ and $136.159 \to 121.016\ms$ &
$2.846\ms$ & $15.143\ms$ & Replaced zero/scatter/readback scratch traffic with parent gathers, shrank topology loads, and used the best launch shape. \\
\bottomrule
\end{longtable}

\section{Instruction-level changes}

\subsection{Leaf species index}

The old path materialized a dense leaf contribution over all wave rows and
species.  The new path carries a compact leaf-species index per clade row and
lets each Triton program decide whether its species lane is the leaf species.

\begin{lstlisting}
# old idea
leaf_term = full((W, S), -inf)
leaf_term[row, leaf_species[row]] = log_pS[leaf_species[row]]
term = combine(term, leaf_term)

# new idea inside the wave kernel
s_leaf = leaf_species_index[row]
leaf_value = where(s == s_leaf, log_pS[s], -inf)
term = combine(term, leaf_value)
\end{lstlisting}

This helped because the representative wave kernels were bandwidth-bound.  The
optimization removed dense $[W,S]$ allocation and read traffic for data that is
mostly \(-\infty\).

\subsection{Read-only RHS and final RHS view}

The critical correctness fact is that the fused wave kernel only reads the
current wave RHS.  Later DTS and \code{Pibar} adjoints write to earlier child
rows, not to the current wave slice.  Therefore the current slice does not need
a clone.

\begin{lstlisting}
# old
rhs_k = accumulated_rhs[ws:we].clone()
wave_backward(rhs_k, ...)

# new
rhs_k = accumulated_rhs[ws:we]       # read-only view
wave_backward(rhs_k, ...)
\end{lstlisting}

The final view change removed 45 D2D copies for 10 families and 49 D2D copies
for 50 families.  At 50 families, D2D copy time fell from $5.159\ms$ to
$0.391\ms$ and D2D bytes fell by $2574.152\mb$.

\subsection{Fused wave parameter accumulation}

The old path wrote multiple full contribution matrices and reduced them later
with PyTorch.  The fused path accumulates the relevant species contribution
inside \code{\_wave\_backward\_uniform\_kernel}.

\begin{lstlisting}
# old
aw0, aw1, aw2, aw3, aw4, aw5 = wave_kernel_outputs(...)
grad_E    += (aw0 + aw2).sum(dim=0)
grad_D    += aw3.sum(dim=0)
grad_S    += aw4.sum(dim=0)
grad_T    += aw5.sum(dim=0)

# new, in the Triton program
local_E = aw0_s + aw2_s
local_D = aw3_s
local_S = aw4_s
local_T = aw5_s
tl.atomic_add(grad_E + s, local_E)
tl.atomic_add(grad_D + s, local_D)
tl.atomic_add(grad_S + s, local_S)
tl.atomic_add(grad_T + s, local_T)
\end{lstlisting}

NCU showed global RED operations after this change, but the tradeoff was
favorable: the largest leaf wave removed roughly $1.82\gb$ of DRAM traffic and
the 50-family wall time improved by $23.919\ms$ in that optimization ladder.

\subsection{Compact \code{Pibar} scratch}

The old Neumann scratch stored several arrays:

\begin{lstlisting}
# old scratch
pibar_wt
inv_denom
p_prime
\end{lstlisting}

The compact path stores the coefficient that is actually needed later and
recomputes the cheap normalized probability when needed:

\begin{lstlisting}
# new scratch
pibar_coeff = pibar_wt * inv_denom

# later
u_d = term * pibar_coeff
p_prime = exp2(Pi - row_max)
\end{lstlisting}

The largest profiled 10-family leaf wave went from $4.032\gb$ to $2.694\gb$
of total DRAM traffic after compact scratch.  Occupancy stayed about the same,
so the win came from issuing fewer memory transactions, not from more resident
work.

\subsection{Kernelized active-mask construction}

The original mask construction used several PyTorch kernels:

\begin{lstlisting}
clade_max = rhs_k.abs().max(dim=1).values
active_mask = clade_max >= pruning_threshold
\end{lstlisting}

The accepted version uses a Triton row-reduction kernel:

\begin{lstlisting}
active_mask = active_mask_from_rhs_absmax(rhs_k, pruning_threshold)
\end{lstlisting}

This does not remove the host pruning decisions.  The Python loop still uses
\code{any()}, \code{sum().item()}, and \code{nonzero()} to decide which waves
and rows to process.  It helped anyway because it replaced about $8.1\ms$ of
PyTorch active-mask GPU work with about $2.95\ms$ of Triton work in the
50-family profile.

\subsection{Merged DTS S-term accumulation}

The old DTS accumulation kernel made two species passes.  The second pass
existed mostly for the S term and reloaded values that had already been touched.

\begin{lstlisting}
# old
for s in species:
    compute D term
    compute T terms
    write grad_Pibar_l/r
    accumulate scalar params

for s in species:
    reload Pi_l, Pi_r, parent Pi, v_k
    compute S term
    atomic_add accumulated_rhs

# new
for s in species:
    compute D term
    compute T terms
    compute S term while masks and values are live
    write grad_Pibar_l/r
    accumulate scalar params and direct Pi adjoints
\end{lstlisting}

NCU confirmed the mechanism: on the largest DTS launch, total DRAM bytes fell
from $3.644\gb$ to $3.345\gb$ and excessive global sectors fell by about
$132\mb$.  The cost was higher register pressure, from 54 to 78 registers per
thread, and lower occupancy.  The net effect was still positive: the DTS
accumulation bucket fell by $3.535\ms$ in Nsight Systems.

\subsection{DTS scalar parameters stay on device}

The old wrappers converted CUDA scalar tensors into Python floats before each
Triton launch:

\begin{lstlisting}
# old
pD_val = float(log_pD.item())
pS_val = float(log_pS.item())
dts_kernel(..., pD_val, pS_val)

# new
dts_kernel(..., log_pD_device_ptr, log_pS_device_ptr)

# inside the kernel
log_pD = tl.load(log_pD_ptr)
log_pS = tl.load(log_pS_ptr)
\end{lstlisting}

For the 50-family workload there are 33 split waves, hence 66 scalar
extractions for \code{log\_pD} and \code{log\_pS}.  Nsight Systems measured
exactly 66 fewer \code{cudaStreamSynchronize} calls, 66 fewer
\code{cudaMemcpyAsync} calls, and 66 fewer D2H copy events.  The DTS kernel
itself did not get faster; the end-to-end interval got faster because the host
no longer stopped to read those scalars.

\subsection{DTS scalar reduction accumulation}

The scalar-parameter path used to materialize per-split arrays and then reduce
them:

\begin{lstlisting}
# old
param_pD, param_pS = dts_backward(...)
grad_log_pD += param_pD.sum()
grad_log_pS += param_pS.sum()
\end{lstlisting}

The promoted path accumulates only the two scalar parameter gradients inside
large split waves:

\begin{lstlisting}
# new inside the DTS accumulation kernel
sum_pD = reduce_over_species(local_pD_contrib)
sum_pS = reduce_over_species(local_pS_contrib)
tl.atomic_add(grad_log_pD, sum_pD)
tl.atomic_add(grad_log_pS, sum_pS)
\end{lstlisting}

This was a small cleanup, not a major kernel optimization.  It removed 36 tiny
PyTorch launches in the 50-family Nsight capture and improved the 50-family
mean by about $0.9\ms$, while the median was essentially neutral.  The
\code{grad\_mt} vector accumulation variant was implemented but left
default-off because it introduced many atomics into only $S=1999$ species
targets and made the register-sensitive DTS kernel slower.

\subsection{DTS-staged \code{Pibar} VJP}

This was the largest split-side improvement.  The old path materialized
\code{grad\_Pibar\_l/r}; the tree VJP then read those matrices, computed
\code{u\_d}, wrote \code{u\_d} into a subtree buffer, and finally performed the
tree correction.

\begin{lstlisting}
# old
grad_Pibar_l, grad_Pibar_r = dts_backward_accum(...)

for side in (left, right):
    inv_denom = compute_denominator(Pi_child, mt)
    u_d = grad_Pibar_side * inv_denom
    subtree_buf = u_d
    pibar_tree_vjp(subtree_buf)
\end{lstlisting}

The new path computes the exact input needed by the tree VJP in the DTS kernel:

\begin{lstlisting}
# new
# DTS already has vd1/vd2, where:
#   grad_Pibar_l = vd2
#   grad_Pibar_r = vd1

inv_denom = exp2(row_max + mt - Pibar_star)
u_d = vd * inv_denom
A = sum_s(u_d[s])

write u_d_staging
write A
pibar_tree_vjp_from_ud(u_d_staging, A)
\end{lstlisting}

The stable identity
\[
  \mathrm{inv\_denom}_{r,s}
    = 2^{\mathrm{row\_max}_r + mt_s - Pibar_{r,s}}
\]
allows the DTS kernel to use forward-stored row maxima and the already
materialized \code{Pibar} value.

The tradeoff was clear in the profiles.  DTS accumulation became slower
($26.237\ms \to 30.006\ms$ in Nsight) because it writes the staged \code{u\_d},
stores row sums, and carries more live values.  But the \code{Pibar} tree bucket
fell from $37.888\ms$ to $25.510\ms$, and PyTorch float reductions fell from
$5.357\ms$ to $0.369\ms$.  NCU showed the tree kernel instructions falling from
$3.085$ billion to $0.900$ billion, registers per thread falling from 56 to 40,
and achieved occupancy rising from $74.58\%$ to $98.37\%$.

\subsection{Wave speciation gather, int32 topology, and row-max reuse}

After the staged \code{Pibar} work, the wave kernel became the largest bucket
again.  Source counters showed scratch traffic around Neumann speciation as the
main issue.  The old code effectively did this:

\begin{lstlisting}
# old scatter-style speciation
zero(spec_buf_next)

for parent in wave_rows:
    spec_buf_next[child1[parent], :] += term[parent, :] * sl1[parent, :]
    spec_buf_next[child2[parent], :] += term[parent, :] * sl2[parent, :]

for child in wave_rows:
    spec = spec_buf_next[child, :]
    next_term[child, :] = combine(spec, ...)
\end{lstlisting}

The new path asks each child row to gather from its parent:

\begin{lstlisting}
# new gather-style speciation
parent = parent_index[child]
side_weight = where(child == child1[parent], sl1[parent, :], sl2[parent, :])
spec = term[parent, :] * side_weight
next_term[child, :] = combine(spec, ...)
\end{lstlisting}

This adds parent-index loads but removes full-row zeroing, scatter stores, and
readback.  The same change also made 32-bit topology useful: species-tree
indices fit in int32, so the kernel avoids repeated 64-bit topology loads.  The
winning launch shape for this gather path was \code{num\_warps=8}; larger
\code{BLOCK\_S} and 16 warps were slower.

Representative NCU numbers for the large wave kernel:

\begin{center}
\begin{tabular}{lrr}
\toprule
Metric & Old scatter & Final default \\
\midrule
Duration & $7.847\ms$ & $5.102\ms$ \\
DRAM read & $2.941\gb$ & $1.008\gb$ \\
DRAM write & $2.437\gb$ & $1.895\gb$ \\
L1 global load bytes & $40.352\gb$ & $20.259\gb$ \\
L1 global store bytes & $6.600\gb$ & $3.518\gb$ \\
L1 global RED bytes & $2.638\gb$ & $1.327\gb$ \\
Achieved occupancy & $82.51\%$ & $99.43\%$ \\
Registers/thread & 48 & 40 \\
Excessive L2 sectors & $777.6\mb$ & $146.1\mb$ \\
\bottomrule
\end{tabular}
\end{center}

The final wave-kernel proposal moved the 50-family median from about
$136.2\ms$ to about $121.0\ms$.  Nsight Systems attributed almost all of the
gain to \code{\_wave\_backward\_uniform\_kernel}, whose bucket fell from
$39.422\ms$ to $25.412\ms$.

\section{Rejected or diagnostic experiments}

Several optimizations were implemented or prototyped but deliberately not
promoted.  They are useful because they explain where not to spend the next
engineering pass.

\begin{longtable}{p{0.24\linewidth}p{0.22\linewidth}p{0.44\linewidth}}
\toprule
Experiment & Result & Interpretation \\
\midrule
\endfirsthead
\toprule
Experiment & Result & Interpretation \\
\midrule
\endhead

Disable pruning & 50-family time $230.795\ms \to 324.348\ms$ in the early table; later $157.329\ms \to 241.924\ms$. & Pruning still skips too much real work to remove it wholesale. \\

Fixed-schedule device pruning & 50-family median about $174.387\ms \to 177.588$--$177.943\ms$ in the first profile; later $157.329\ms \to 161.084$--$161.736\ms$. & It removed D2H syncs, but launched kernels for waves the host-pruned path skipped. No-op scheduling cost exceeded sync savings. \\

Pruning threshold $10^{-8}$ & $230.795\ms \to 218.066\ms$ from the early baseline, but slower than the accepted path. & More aggressive thresholding was not the right source of speedup and risked changing pruning behavior. \\

Child-grouped uniform \code{Pibar} VJP & Tree VJP bucket $37.784\ms \to 32.119\ms$, but group-reduce cost was $16.102\ms$; net path $48.221\ms$. & Duplicate split-side factor was only about $1.111\times$, too small to pay for a separate reduction pass. \\

Forward \code{Pibar} stat reuse alone & 10-family median $47.769\ms \to 47.290\ms$, but 50-family median $162.704\ms \to 166.956\ms$. & It reduced ancestor-gather work but added global \code{Pibar}/\code{mt} reads; at 50 families the DRAM pressure won. \\

Backward row-stat precompute & Tree VJP improved by $0.848\ms$, but a $2.690\ms$ precompute pass was added. & A full extra pass over row data was more expensive than the saved denominator work. \\

Species-tree prefix denominator & 50-family median $162.704\ms \to 180.791\ms$; Pibar VJP bucket $37.784\ms \to 56.475\ms$. & It reduced some gather stalls but introduced level barriers and more L2/DRAM traffic. \\

Grouped DTS backward accumulation & 50-family median $164.352\ms \to 175.860\ms$. & High-fanout waves were not child-duplicate-heavy; grouping added materialization and reductions. \\

No-atomic DTS accumulation & 50-family median $164.352\ms \to 163.053\ms$, but Nsys DTS bucket was unchanged. & Correct for child-unique cases, but atomics were not the dominant cost in this layout. \\

Disable kernelized backward DTS & 50-family time $180.732\ms \to 649.828\ms$ in the early comparison. & The fused DTS backward path is essential. \\

No-split final VJP shortcut & 50-family $181$--$184\ms$ range, slower than compact scratch. & Algebra was valid, but it read scratch and increased DRAM traffic. \\

Recompute \code{Pibar} denominator scratch & 50-family $174.754\ms$, slightly worse than compact coefficient scratch $174.174\ms$. & Diagnostic useful for understanding traffic, but compact scratch was better. \\

Leaf-logp hit-only/scalar specialization & 50-family $182.300\ms$, worse than the accepted path. & Leaf log-probability loads were not the limiting traffic after leaf indexing. \\

Block/warp sweep for wave kernel & \code{BLOCK\_S=128}: $143.905\ms$; \code{BLOCK\_S=512}: $137.661\ms$; 16 warps: $125.623\ms$. & The winning point was \code{BLOCK\_S=256}, 8 warps on the gather path. Larger shapes increased loop trips, pressure, or locality cost. \\

\bottomrule
\end{longtable}

\section{Why the accepted changes worked}

The successful changes fall into four categories.

\subsection{Fewer bytes}

This was the dominant theme.  Leaf indexing, compact \code{Pibar} scratch, RHS
view mode, staged \code{Pibar} VJP, and wave speciation gather all removed large
global-memory streams.  The clearest examples are:

\begin{itemize}[leftmargin=*]
\item compact \code{Pibar} scratch removed about $1.34\gb$ of DRAM traffic from
      one representative 10-family leaf-wave launch;
\item RHS view removed $2574.152\mb$ of D2D copies at 50 families;
\item staged \code{Pibar} VJP reduced representative tree-kernel DRAM read by
      about $681\mb$ and cut executed instructions by about $71\%$;
\item wave gather reduced representative wave-kernel DRAM read by about
      $1.93\gb$ and excessive L2 sectors by about $81\%$.
\end{itemize}

\subsection{Fewer launches and PyTorch reductions}

Fused wave parameter accumulation, direct global reductions, DTS scalar
reduction accumulation, and staged \code{Pibar} VJP removed PyTorch reduction
kernels from the hot path.  Some of these were small individually, but they
also removed large temporary tensors.  Proposal 5, for example, reduced total
kernel launches from 3100 to 2968 in its Nsight capture and cut PyTorch
float-sum reductions from $5.357\ms$ to $0.369\ms$.

\subsection{Fewer host synchronizations}

Device-resident DTS scalars did not change the bulk GPU kernel profile.  Its
value was removing 66 scalar D2H reads and their stream synchronizations.
Kernelized active-mask construction did not remove the host pruning decisions,
but it did remove several PyTorch kernels used to build the mask.

\subsection{Better data layout}

The wave gather path and int32 topology change are layout optimizations.  The
important change was not more arithmetic parallelism; it was changing the
memory access pattern so the kernel no longer writes a full speciation scratch
buffer only to read it back.  Int32 topology then reduced the width of repeated
index loads.

\section{Correctness discipline}

The reports consistently paired performance work with parity checks:

\begin{itemize}[leftmargin=*]
\item old-vs-new loss and theta-gradient comparisons on 3, 10, and 50 family
      workloads;
\item focused Triton kernel tests for active mask, DTS backward accumulation,
      uniform cross-\code{Pibar} VJP, staged DTS+\code{Pibar}, and scatter-vs-
      gather wave behavior;
\item finite-difference/autograd bridge tests, including forced fallback paths
      through environment variables;
\item fp32/fp64 direct parity for the DTS scalar, reduction, and staged
      \code{Pibar} paths where applicable.
\end{itemize}

Small fp32 differences were accepted only where the summation order changed,
for example the fused parameter accumulation path.  The documented maximum
relative differences were on the order expected from reordered fp32 reductions,
and the finite-difference bridge tests passed.

\section{Current bottleneck after these changes}

After proposal 6, the wave kernel is no longer the dominant bucket.  The next
large costs are again split-side work: DTS accumulation and the staged
\code{Pibar} tree VJP.  The successful changes also make clear what future work
should avoid.  A monolithic DTS+\code{Pibar} kernel is unattractive because the
staged DTS kernel already rose to 96 registers per thread.  Device pruning must
preserve the useful scheduling property of host pruning; simply launching all
waves with active masks was slower.

The next likely direction is a device-resident active worklist or a more
parent-aware DTS layout: preserve pruning's ability to skip work, avoid host
scalar decisions, and reduce the high-fanout split-wave traffic without adding
large materialized reduction passes.

\end{document}
